\documentclass[10pt,a4paper]{amsart}
\usepackage{amsmath,amssymb,amsthm,graphicx}
% Retained minimal preamble for the Stacks Project chapter more-algebra.tex (More on Algebra).
% The chapter file pinned-original-source.tex carries only \input{preamble} on its first line, so no
% preamble text lives inside the chapter itself.  The Stacks Project uses one single shared preamble for
% every chapter; the snapshot stages it identically as sources/stacks-chow/preamble.tex and
% sources/stacks-etale-cohomology/preamble.tex.  The block below is copied mechanically, line for line,
% from sources/stacks-chow/preamble.tex lines 14-18 (the Xy-pic setup) and 154-217 (the native theorem
% declarations together with the mathematical macros).  Nothing was added, removed or redefined; the
% wrapper class and the amsmath/amssymb/amsthm/graphicx base packages are supplied by the assembler.
% Chapter: more-algebra.tex -- More on Algebra.
% For commutative diagrams we use Xy-pic
\usepackage[all]{xy}

% We use 2cell for 2-commutative diagrams.
\xyoption{2cell}

% Theorem environments.
%
\theoremstyle{plain}
\newtheorem{theorem}[subsection]{Theorem}
\newtheorem{proposition}[subsection]{Proposition}
\newtheorem{lemma}[subsection]{Lemma}

\theoremstyle{definition}
\newtheorem{definition}[subsection]{Definition}
\newtheorem{example}[subsection]{Example}
\newtheorem{exercise}[subsection]{Exercise}
\newtheorem{situation}[subsection]{Situation}

\theoremstyle{remark}
\newtheorem{remark}[subsection]{Remark}
\newtheorem{remarks}[subsection]{Remarks}

\numberwithin{equation}{subsection}

% Macros
%
\def\lim{\mathop{\mathrm{lim}}\nolimits}
\def\colim{\mathop{\mathrm{colim}}\nolimits}
\def\Spec{\mathop{\mathrm{Spec}}}
\def\Hom{\mathop{\mathrm{Hom}}\nolimits}
\def\Ext{\mathop{\mathrm{Ext}}\nolimits}
\def\SheafHom{\mathop{\mathcal{H}\!\mathit{om}}\nolimits}
\def\SheafExt{\mathop{\mathcal{E}\!\mathit{xt}}\nolimits}
\def\Sch{\mathit{Sch}}
\def\Mor{\mathop{\mathrm{Mor}}\nolimits}
\def\Ob{\mathop{\mathrm{Ob}}\nolimits}
\def\Sh{\mathop{\mathit{Sh}}\nolimits}
\def\NL{\mathop{N\!L}\nolimits}
\def\CH{\mathop{\mathrm{CH}}\nolimits}
\def\proetale{{pro\text{-}\acute{e}tale}}
\def\etale{{\acute{e}tale}}
\def\QCoh{\mathit{QCoh}}
\def\Ker{\mathop{\mathrm{Ker}}}
\def\Im{\mathop{\mathrm{Im}}}
\def\Coker{\mathop{\mathrm{Coker}}}
\def\Coim{\mathop{\mathrm{Coim}}}

% Boxtimes
%
\DeclareMathSymbol{\boxtimes}{\mathbin}{AMSa}{"02}

%
% Macros for moduli stacks/spaces
%
\def\QCohstack{\mathcal{QC}\!\mathit{oh}}
\def\Cohstack{\mathcal{C}\!\mathit{oh}}
\def\Spacesstack{\mathcal{S}\!\mathit{paces}}
\def\Quotfunctor{\mathrm{Quot}}
\def\Hilbfunctor{\mathrm{Hilb}}
\def\Curvesstack{\mathcal{C}\!\mathit{urves}}
\def\Polarizedstack{\mathcal{P}\!\mathit{olarized}}
\def\Complexesstack{\mathcal{C}\!\mathit{omplexes}}
% \Pic is the operator that assigns to X its picard group, usage \Pic(X)
% \Picardstack_{X/B} denotes the Picard stack of X over B
% \Picardfunctor_{X/B} denotes the Picard functor of X over B
\def\Pic{\mathop{\mathrm{Pic}}\nolimits}
\def\Picardstack{\mathcal{P}\!\mathit{ic}}
\def\Picardfunctor{\mathrm{Pic}}
\def\Deformationcategory{\mathcal{D}\!\mathit{ef}}

% --- appended 2026-10-12 (Chief): the project own definitions of the reference/slogan/history
% display environments, copied mechanically from the shared Stacks preamble lines 8-12
% (\usepackage{verbatim} plus the three \newenvironment lines, which the chapter preamble keeps
% verbatim from the project preamble). Nothing added, removed or redefined: these environments
% carry no mathematical content and the project itself typesets them as comments. Without these
% three lines a counted statement whose own text uses them fails with Environment slogan undefined.
\usepackage{verbatim}
\newenvironment{reference}{\comment}{\endcomment}
\newenvironment{slogan}{\comment}{\endcomment}
\newenvironment{history}{\comment}{\endcomment}
\usepackage{cleveref}
\crefname{originallemma}{original lemma}{original lemmas}
\Crefname{originallemma}{Original lemma}{Original lemmas}
\makeatletter
\newlabel{lemma-cone-relatively-perfect}{{0DHT}{source}{}{sourcepointer-1}{}}
\newlabel{lemma-cone-relatively-perfect@cref}{{[originallemma][0][]0DHT}{[1][0][]source}{}{}{}}
\makeatother
\crefname{originallemma}{original lemma}{original lemmas}
\Crefname{originallemma}{Original lemma}{Original lemmas}
\makeatletter
\newlabel{lemma-Noetherian-pseudo-coherent}{{066E}{source}{}{sourcepointer-2}{}}
\newlabel{lemma-Noetherian-pseudo-coherent@cref}{{[originallemma][0][]066E}{[1][0][]source}{}{}{}}
\makeatother
\crefname{originallemma}{original lemma}{original lemmas}
\Crefname{originallemma}{Original lemma}{Original lemmas}
\makeatletter
\newlabel{lemma-base-change-comparison}{{0661}{source}{}{sourcepointer-3}{}}
\newlabel{lemma-base-change-comparison@cref}{{[originallemma][0][]0661}{[1][0][]source}{}{}{}}
\makeatother
\crefname{originallemma}{original lemma}{original lemmas}
\Crefname{originallemma}{Original lemma}{Original lemmas}
\makeatletter
\newlabel{lemma-last-one-flat}{{0653}{source}{}{sourcepointer-4}{}}
\newlabel{lemma-last-one-flat@cref}{{[originallemma][0][]0653}{[1][0][]source}{}{}{}}
\makeatother
\crefname{originallemma}{original lemma}{original lemmas}
\Crefname{originallemma}{Original lemma}{Original lemmas}
\makeatletter
\newlabel{algebra-lemma-flat-finite-presentation-limit-flat}{{02JO}{source}{}{sourcepointer-5}{}}
\newlabel{algebra-lemma-flat-finite-presentation-limit-flat@cref}{{[originallemma][0][]02JO}{[1][0][]source}{}{}{}}
\makeatother
\title{An object is R-perfect exactly when it is a finite complex of flat finitely presented modules}
\author{The Stacks Project Authors}
\date{stacks-project:more-algebra.tex:local-snapshot-20261012; GNU FDL 1.2 or later}
\begin{document}
\maketitle
\setcounter{subsection}{1}
\noindent\textit{Editorial note.} All delivered mathematical text is the author's own, copied from the pinned Stacks Project chapter file more-algebra.tex; no statement or proof step was rewritten, reordered or removed. The counted span is the author environment and its own proof environment. References inside the span are handed to the assembler as navigation stubs that carry the canonical tag of the cited tag only; whether each of them is load-bearing is recorded on the stub itself.\par\medskip
\noindent\textit{Allowed standard prerequisites (editorial).}\begin{itemize}
\item An object of finite tor amplitude over the base is R-perfect as soon as it is pseudo-coherent. used at proof lines 21990-21995. Reduces the implication from (2) to (1) to showing that a flat finitely presented module is pseudo-coherent.
\item A flat finitely presented module over a base descends to a finitely generated situation, where finite type modules over Noetherian rings are pseudo-coherent. invoked at proof lines 21995-22010. Produces the finitely generated model M\_0 over a finitely generated base and its pseudo-coherence, and the finite free resolution that spreads out.
\item A bounded above complex of finite free modules whose truncation cuts at the tor amplitude lower bound consists of flat finitely presented modules. invoked at proof lines 22015-22019. Supplies the representing complex for the implication from (1) to (2).
\end{itemize}\medskip
\ifdefined\hypertarget\hypertarget{sourcepointer-1}{}\fi
\noindent\textit{Source reference.} Original lemma 0DHT: sources/stacks-more-algebra/pinned-original-source.tex, source lines 21949--21949. The sentence at proof lines 21990-21992 invokes it to reduce (2) implies (1) to a single R-flat finitely presented module; the verification of that case is carried out inside the counted span.\par\medskip
\ifdefined\hypertarget\hypertarget{sourcepointer-2}{}\fi
\noindent\textit{Source reference.} Original lemma 066E: sources/stacks-more-algebra/pinned-original-source.tex, source lines 16517--16517. Used at proof line 22003 to conclude that the descended module M\_0 is pseudo-coherent; the spreading out of M\_0, the finite free resolution and the base change back to A are inside the counted span.\par\medskip
\ifdefined\hypertarget\hypertarget{sourcepointer-3}{}\fi
\noindent\textit{Source reference.} Original lemma 0661: sources/stacks-more-algebra/pinned-original-source.tex, source lines 14955--14955. Cited parenthetically at proof lines 22008-22009 to see that the base changed finite free resolution still resolves M; the flatness of A\_0 over R\_0 and the flatness of M\_0 over R\_0 that make that comparison available are used in the counted span.\par\medskip
\ifdefined\hypertarget\hypertarget{sourcepointer-4}{}\fi
\noindent\textit{Source reference.} Original lemma 0653: sources/stacks-more-algebra/pinned-original-source.tex, source lines 16883--16883. Used at proof line 22018 for the implication from (1) to (2): it says the truncated finite free resolution consists of flat finitely presented modules; the truncation step itself and the identification of the representing complex are inside the counted span.\par\medskip
\ifdefined\hypertarget\hypertarget{sourcepointer-5}{}\fi
\noindent\textit{Source reference.} Original lemma 02JO: sources/stacks-chow/tags.txt, source lines 1--1. The reference at proof line 21995 produces the finitely generated model R\_0 to A\_0 and the module M\_0; everything done with that model afterwards, including the resolution and the base change, is inside the counted span.\par\medskip
\begin{lemma}
\label{lemma-structure-relatively-perfect}
Let $R \to A$ be a flat ring map of finite presentation.
Let $K \in D(A)$. The following are equivalent
\begin{enumerate}
\item $K$ is $R$-perfect, and
\item $K$ is isomorphic to a finite complex of $R$-flat,
finitely presented $A$-modules.
\end{enumerate}
\end{lemma}

\begin{proof}
To prove (2) implies (1) it suffices by
Lemma \ref{lemma-cone-relatively-perfect}
to show that an $R$-flat, finitely presented $A$-module $M$ defines
an $R$-perfect object of $D(A)$. Since $M$ has finite tor dimension
over $R$, it suffices to show that $M$ is pseudo-coherent. By
Algebra, Lemma \ref{algebra-lemma-flat-finite-presentation-limit-flat}
there exists a finite type $\mathbf{Z}$-algebra $R_0 \subset R$
and a flat finite type ring map $R_0 \to A_0$ and
a finite $A_0$-module $M_0$ flat over $R_0$ such that
$A = A_0 \otimes_{R_0} R$ and $M = M_0 \otimes_{R_0} R$. By
Lemma \ref{lemma-Noetherian-pseudo-coherent}
we see that $M_0$ is pseudo-coherent $A_0$-module.
Choose a resolution $P_0^\bullet \to M_0$ by finite free
$A_0$-modules $P_0^n$. Since $A_0$ is flat over $R_0$,
this is a flat resolution. Since $M_0$ is flat over $R_0$
we find that $P^\bullet = P_0^\bullet \otimes_{R_0} R$
still resolves $M = M_0 \otimes_{R_0} R$. (You can use
Lemma \ref{lemma-base-change-comparison} to see this.)
Hence $P^\bullet$ is a finite free resolution of $M$
over $A$ and we conclude that $M$ is pseudo-coherent.

\medskip\noindent
Assume (1). We can represent $K$ by a bounded above complex
$P^\bullet$ of finite free $A$-modules. Assume that
$K$ viewed as an object of $D(R)$ has tor amplitude in $[a, b]$.
By Lemma \ref{lemma-last-one-flat} we see that
$\tau_{\geq a}P^\bullet$ is a complex of $R$-flat, finitely
presented $A$-modules representing $K$.
\end{proof}

\end{document}
